% GENEROVANO z ../skripta/preamble.tex -- barvy a matematicka makra (sdili TikZ obrazky i kontrola matematiky)
% ---------- Barvy ----------
\definecolor{defc}{HTML}{1F5FA8}   % definice -- modra
\definecolor{thmc}{HTML}{8E2C6E}   % vety / lemmata -- purpurova
\definecolor{exc}{HTML}{2E7D32}    % priklady -- zelena
\definecolor{intc}{HTML}{C77700}   % intuice -- jantarova
\definecolor{warnc}{HTML}{C62828}  % pozor -- cervena
\definecolor{sumc}{HTML}{455A64}   % shrnuti -- sedomodra
\definecolor{cvc}{HTML}{00796B}    % cviceni -- tyrkysova
\definecolor{solc}{HTML}{5D4037}   % reseni -- hneda
\definecolor{acc}{HTML}{2E7D32}    % "prijato" (qA, ANO)
\definecolor{rej}{HTML}{C62828}    % "zamitnuto" (qR, NE)
\definecolor{hl}{HTML}{FFE082}     % zvyrazneni bunky / vrcholu
\definecolor{tT}{HTML}{43A047}     % barva T (True)
\definecolor{tF}{HTML}{E53935}     % barva F (False)
\definecolor{tB}{HTML}{1E88E5}     % barva B (Base)
\definecolor{gr}{HTML}{9E9E9E}

% ---------- Matematicke makra ----------
\newcommand{\cls}[1]{\ensuremath{\mathsf{#1}}}
\newcommand{\PP}{\cls{P}}
\newcommand{\NP}{\cls{NP}}
\newcommand{\NPhat}{\ensuremath{\widehat{\mathsf{NP}}}}
\newcommand{\EXP}{\cls{EXP}}
\newcommand{\NEXP}{\cls{NEXP}}
\newcommand{\DTIME}{\cls{DTIME}}
\newcommand{\NTIME}{\cls{NTIME}}
\newcommand{\DSPACE}{\cls{DSPACE}}
\newcommand{\NSPACE}{\cls{NSPACE}}
\newcommand{\PSPACE}{\cls{PSPACE}}
\newcommand{\NPSPACE}{\cls{NPSPACE}}
\newcommand{\LL}{\cls{L}}
\newcommand{\NL}{\cls{NL}}
\newcommand{\coNP}{\cls{coNP}}
\newcommand{\coNL}{\cls{coNL}}
\newcommand{\redp}{\le_{p}}
\newcommand{\enc}[1]{\ensuremath{\langle #1\rangle}}
\newcommand{\prob}[1]{\textsc{#1}}
\newcommand{\qA}{\ensuremath{q_{A}}}
\newcommand{\qR}{\ensuremath{q_{R}}}
\newcommand{\qhalt}{\ensuremath{q_{\mathrm{halt}}}}
\newcommand{\blank}{\ensuremath{\sqcup}}
\newcommand{\Sig}{\ensuremath{\Sigma}}
\newcommand{\N}{\ensuremath{\mathbb{N}}}
\newcommand{\Lft}{\ensuremath{\leftarrow}}
\newcommand{\Rgt}{\ensuremath{\rightarrow}}
\newcommand{\Stay}{\ensuremath{\bullet}}
\newcommand{\ANO}{\textcolor{acc}{\textsf{\textbf{ANO}}}}
\newcommand{\NE}{\textcolor{rej}{\textsf{\textbf{NE}}}}
\DeclareMathOperator{\dist}{dist}
\DeclareMathOperator{\polylog}{polylog}

\newcommand{\clr}[2]{{\color{#1}#2}}
\providecommand{\uv}[1]{\quotedblbase #1\textquotedblleft}
